\subsection{Overview}

Phosphorus (P) moves in far smaller mass than nitrogen, yet even at low
concentrations it drives eutrophication, and the long-standing presumption
that the soil profile filters subsurface water into a ``clean'' P pathway
has been overturned. A geographically diverse body of monitoring, isotopic
and radionuclide tracing, soil micromorphology, dye and smoke-injection
experiments, and process modelling now supports a well-consolidated
picture. Subsurface drainage typically captures the largest share of field
water and, despite lower P concentrations than surface runoff, frequently
dominates edge-of-field P \emph{loads}, often as particulate P (PP)
delivered by preferential flow
\parencite{Pease2018,VanEsbroeck2016,Uusitalo2001,Laubel1999}. Dissolved P
release is tightly coupled to soil P status, degree of P saturation, and
subsoil sorption, with legacy P sustaining loads at watershed scale
\parencite{Sharpley1995,Beauchemin2003,Jarvie2017}. A tiny, hydraulically
active macropore volume---earthworm burrows, cracks, root channels, and
installation scars---governs rapid solute and sediment delivery and
reorganises soil structure over years to decades
\parencite{Shipitalo2004,Villholth1998,Montagne2008}. Losses are highly
concentrated in time, dominated by the non-growing season and a few large
events \parencite{VanEsbroeck2016,Brendel2019}. This section synthesises
these established magnitudes and mechanisms; P-removal technologies and the
P responses of controlled drainage are treated in the edge-of-field and
design sections, respectively.

\subsection{Phosphorus Loss Pathways and Water Partitioning}

Tile drainage commonly captures the largest share of field water export:
subsurface discharge accounted for $81\pm23\%$ of annual discharge in the
Eastern Corn Belt \parencite{Pease2018} and $78$--$90\%$ in Ontario
\parencite{VanEsbroeck2016}, and installing drainage re-routes water from
surface- to drainage-dominated flow \parencite{Turtola1995,BENGTSON1995}. A
recurring distinction is that surface runoff usually carries higher P
\emph{concentrations} whereas tile drainage---carrying more water---can
dominate annual \emph{loads}. In Ohio, flow-weighted DRP and TP
concentrations were greater in surface runoff, yet subsurface discharge
supplied $71\%$ of DRP and $69\%$ of TP load \parencite{Pease2018}; the same
contrast appears in long-term Swedish monitoring \parencite{Norberg2022}.
Tiles exported $40$--$77\%$ of TP and $19$--$67\%$ of DRP load in Ontario
\parencite{VanEsbroeck2016}, $21$--$52\%$ of TP in Wisconsin
\parencite{Madison2014}, and $65.4\%$ of annual TP at a southern-Ontario
field \parencite{Qi2018}. Models corroborate the partitioning, attributing
$66\%$ of DRP and $76\%$ of TP loss to subsurface drainage against observed
$70\%$ and $80\%$ \parencite{Grewal2025,Sadhukhan2019b}.

Tile dominance is not universal. In snowmelt-affected clay landscapes of the
northern Great Plains, overland flow contributed $72$--$89\%$ of runoff and
$>95\%$ of SRP and TP because tiles were decoupled from the surface by
frozen ground \parencite{Kokulan2019b,Kokulan2021}, and surface runoff also
prevailed under some tillage, pasture, and flooded-paddy systems
\parencite{Zhao2001,Melland2008,DarziNaftchali2016}.

\subsection{Particulate versus Dissolved Phosphorus and Bioavailability}

A central finding is that PP is frequently the dominant form in tile
drainage, contrary to the earlier presumption that the profile filters
subsurface water to dissolved-dominated. In Finnish clays PP was the
grand-mean $92\%$ of TP in both runoff and drainage \parencite{Uusitalo2001};
at Rothamsted's Broadbalk drains PP averaged $54\%$ of TP and was strongly
related to particulate matter \parencite{Lu2003,Lu2004}; and PP dominated
catchment losses in Denmark and elsewhere
\parencite{Grant1996,Brendel2019,vanderSalm2012}. Dissolved forms
nonetheless dominate in snowmelt-driven, grassland, high-STP, and
organic-soil systems: tile SRP averaged $60\%$ of TP across Michigan sites
\parencite{Clement2017}, and DRP dominated fertilised continuous-grass
losses \parencite{Zhang2015b}. Scottish work showed molybdate-reactive P was
only $42\%$ of subsurface TP loss, underscoring the need to measure all P
forms \parencite{Hooda1999}.

Because algal availability, not total mass, governs eutrophication risk,
bioavailable-P estimates matter: desorbable PP and DRP were of equal
importance to algal-available P in Finnish clays
\parencite{Uusitalo2001,Uusitalo2003}, and only $8$--$18\%$ of grassland
lysimeter TP was molybdate-reactive \parencite{Bilotta2008}.

\subsection{Soil Controls on Phosphorus Release and Legacy Phosphorus}

Dissolved P loss is tightly coupled to soil P status. Loss in surface runoff
and drainage related strongly to soil P content
($r^2 = 0.85$--$0.96$ for surface, $0.86$ for subsurface)
\parencite{Sharpley1995}, and higher soil-test P (STP) raised DRP in both
tiles and runoff \parencite{Osterholz2020}. Threshold behaviour is common: a
change point in STP marks where extractable P and DRP rise more steeply
\parencite{McDowell2001b}, and rainfall sharply increased subsurface DRP
where STP exceeded $41$~mg~kg$^{-1}$ \parencite{Ni2020}. Where soils sorb P
weakly, dissolved losses rise: A-horizon degree of P saturation (DPS) alone
explained $83\%$ of tile TP variance, rising to $95$--$96\%$ with
water-extractable P and sorption grouping \parencite{Beauchemin2003}, and a
DPS threshold near $0.25$ predicted sandy-soil SRP \parencite{Siemens2004}.
The highest-leaching profiles combined large labile topsoil P pools with
low-sorption subsoils, so subsoil characteristics must be included in risk
assessment \parencite{Liu2024b}. Subsurface P losses are generally small in
fine-textured soils at moderate STP but become significant in sandy,
organic, and very high-P soils from long-term fertilisation or manuring
\parencite{Sims1998}. Sequential fractionation links pools to loss, with P
lost to water associated with short-term ($<24$~h) exchangeable pools tied to
Al, Fe, and Ca \parencite{McDowell2001}.

Legacy P is manifest at watershed scale: cumulative SRP loads in Lake Erie
tributaries rose $94$--$253\%$ after 2002 even as cropland P balances were
drawn down, $63$--$66\%$ of the increase attributable to greater delivery
consistent with increased tile drainage and connectivity
\parencite{Jarvie2017}.

Organic and drained peat soils are distinctive high-risk sources, where
wetting/drying, organic-matter degradation, and redox chemistry drive
transport \parencite{Grenon2021}. Flooding organic soils increased P release
$4$--$8$-fold, and organic-P mineralisation can supply up to
$160$~kg~P~ha$^{-1}$~yr$^{-1}$ \parencite{Sharpley1995}. Very large filterable
P loads were measured from an Organic soil ($86.7$~kg~P~ha$^{-1}$) versus a
Podzol ($1.7$) \parencite{McDowell2015b}, and degraded drained peatlands were
the highest-risk eutrophication source in a German catchment, governed by
seasonal, redox-mediated remobilisation \parencite{Tiemeyer2009,YliHalla2022}.
Storm-triggered ochre flocs from tile drains carried $27$--$141$~mg~L$^{-1}$
TP with Fe and P co-distributed \parencite{Zimmer2016}.

\subsection{Preferential Flow, Macroporosity, and Drainage-Induced Soil Structure}

\subsubsection{Biopores as the dominant conduits}

Rapid delivery of water and surface-applied solutes to drains is governed by
a small volume of continuous macropores rather than uniform matrix flow. In a
Finnish clay, at $10$~cm above the tile movement occurred exclusively through
earthworm burrows ($106$ dyed burrows~m$^{-2}$), the base of the plough layer
acted as a hydraulic barrier, and infiltration above the tile lines was about
twice that between them and correlated with \textit{Lumbricus terrestris}
numbers \parencite{Shipitalo2004}. Worm burrows were the only structures at
tile depth in a Swiss loam \parencite{Stamm2002}, and roughly $50$ continuous
burrows~m$^{-2}$ reached from surface to $90$~cm at a German site, their
abundance the most important precondition for preferential flow of isoproturon
\parencite{Zehe2001}. Smoke-injection and dye studies confirmed continuous
surface-to-drain pathways clustered within a $\sim$1-m band over the drain,
their number increasing as the water table dropped
\parencite{Nielsen2015,Petersen2012}. The hydraulically active pore volume is
tiny: field analysis indicated an active macroporosity of only $\sim0.17\%$ of
soil volume, yet this fraction produced abrupt drainage concentration
increases by factors of $20$--$400$ within hours \parencite{Villholth1998},
and macropore area fractions peaked below $0.009$, correlating strongly with
saturated conductivity \parencite{Frey2011}.

\subsubsection{Spatial organisation and static structure}

Two organising gradients recur: macroporosity and biopore abundance decline
sharply with depth, whereas lateral variation relative to the drain is real
but weaker. In a Finnish clay, macroporosity fell strongly with depth but
differed little with distance, though earthworm burrows were $44\%$ more
numerous above a drain and field saturated conductivity varied with drain
position \parencite{Alakukku2010,Frey2011}. Reviews report higher
infiltration and lower bulk density above drain lines and runoff reductions of
$34$--$55\%$ relative to undrained soils \parencite{Montagne2009}. Sustained
drainage consistently lowers bulk density and raises air-filled porosity and
gaseous diffusivity: Ohio drainage decreased bulk density and increased
air-filled pores \parencite{HUNDAL1976}, and relative gas-diffusion gains
correlated with corn yield \parencite{Nakajima2014}. Consistent with an
aeration mechanism, drainage raised soil oxygen, mineralisation, and
nitrification \parencite{GROFFMAN1984}, and dynamic aeration indicators
outperformed static bulk density in explaining productivity across a large
Canadian survey \parencite{Caron2025}. The effect is context-dependent,
however: some comparisons found lower bulk density under undrained or
subirrigated soils \parencite{LAL1993,Baker2004}, drainage reduced surface
organic matter and soil organic carbon in some Ohio studies
\parencite{FAUSEY1992,Kumar2014}, and in the most compacted soils aeration
control by drainage becomes impossible \parencite{SMEDEMA1993}.

\subsubsection{Hydraulic evolution, mole drainage, and pedogenesis}

Drainage and associated loosening raise infiltration and saturated hydraulic
conductivity, and these gains often develop over years. Fields drained
$>15$~years earlier had saturated conductivity $7.58$~cm~h$^{-1}$ faster than
recently drained fields, with preferential-flow development the likely driver,
implying that soil physical and biological properties evolve after
installation \parencite{Frankl2023,JafariTalukolaee2016}. Where matrix
permeability is inadequate, mole drainage and subsoiling create the effective
pore network: up to $90\%$ of flow occurred through mole drains on UK heavy
clays \parencite{HARRIS1993}, and moled structure persists as open cavities
for $3$--$7$ years and thereafter as zones of increased permeability
\parencite{MULLER1988,PLAMENAC1988,Castanheira2010}. Over years to decades,
redirected water fluxes accelerate pedogenesis: subsurface drainage
intensified eluviation and clay translocation, removed swelling clays and
$>$one-third of Fe oxides near the drain, and redistributed Mn oxides, with
eluviation restricted to the first metre either side of the drain and
redox-driven changes extending several metres
\parencite{Montagne2010,Montagne2008,Montagne2013,Mercier2000}. About
$70$~years of drainage formed an argic horizon in boreal Finland, a change
judged probably partly irreversible \parencite{YliHalla2009}.

\subsubsection{Macropore transport of solutes and manure constituents}

These structural findings manifest as rapid, connectivity-limited preferential
transport. Conservative tracers reach drains within minutes to hours: chloride
breakthrough increased by factors of $20$--$400$ within hours
\parencite{Villholth1998}; bromide appeared within $1$~h yet cumulative
capture was slow ($8\%$ in $48$~h, $\sim95\%$ only after months)
\parencite{Frey2012,Kladivko2023}. Liquid-manure constituents reached drains
within $5$~min, and soluble P moved through artificial macropores with little
sorption, identifying macropore flow as the primary rapid-P mechanism
\parencite{Stamm2002,Geohring2001}. These behaviours are captured only by
models representing two-domain flow or a direct macropore--drain connection: a
two-dimensional mobile--immobile model matched bromide dynamics far better than
single-porosity formulations \parencite{Köhne2006}; adding a macropore domain
was essential to match drainage chloride \parencite{Villholth1998b}; an
``express fraction'' of $\sim2\%$ reproduced immediate tracer breakthrough in
RZWQM \parencite{Fox2004,Fox2007b}; and a macropore subroutine in APEX raised
DRP-load performance from unsatisfactory to very good, with median annual
macropore-flow and macropore-P contributions of $48\%$ and $43\%$
\parencite{Ford2017}. Fine-textured soils generated more macropore flow than
coarse-textured soils, so soil texture explained connectivity-driven seasonal
variation in flow partitioning \parencite{Poon2021,Holbak2021}.

\subsection{Temporal Dynamics of Phosphorus and Sediment Loss}

Phosphorus loss is highly concentrated in time. Most annual loss occurs in the
non-growing season, which accounted for $83$--$97\%$ of runoff, $84$--$100\%$
of DRP, and $67$--$98\%$ of TP loss in Ontario
\parencite{VanEsbroeck2016,Qi2018}. Event dominance is stark: more than $75\%$
of DRP and $66\%$ of TP export occurred in the top $25\%$ of flows, and a
single 6-day event supplied $46.6\%$ of TP, $84.0\%$ of DRP, and $81.0\%$ of
suspended solids in one subwatershed \parencite{Brendel2019}.
Concentration--discharge behaviour differs by pathway: tile TP and DRP
increased with discharge (chemodynamic) while surface runoff was more
chemostatic, and the tile DRP:TP ratio fell at high flow, indicating a growing
PP contribution \parencite{Osterholz2020}. In cold regions, snowmelt and
winter thaws dominate exports, predominantly in dissolved form---limiting the
value of particulate-focused practices \parencite{Baulch2019,Clement2017}. A
drought-then-deluge sequence in 2021 produced unprecedented P export ($43\%$
of a 12-year DRP load in nine months) \parencite{Millar2025,Mumbi2025b}.

Preferential flow reconciles substantial subsurface PP loss with the filtering
expected of matrix flow. $^{137}$Cs tracing shows drainage particulate matter
originates from topsoil (median grain sizes $2$--$5$~$\mu$m) translocated
within hours \parencite{Laubel1999,Uusitalo2001}, and PP loss is
episodic---nearly one-third of annual PP loss occurred during sampled storms
\parencite{Laubel1999,Grant1996}. Because of this episodicity, discrete
sampling underestimated annual P loss by $26\%$ to $>100\%$
\parencite{Grant1996}. At catchment scale, source-fingerprinting attributed
$28$--$55\%$ of sediment yield to field drains, with topsoil supplying
$70.9$--$88.6\%$ of the load and drains contributing over $50\%$ of the annual
sediment budget in one catchment \parencite{Walling2002,Chapman2005}.

\subsection{In-Field Controls on Phosphorus Loss}

Long-term P inputs raise subsurface P losses, generally in the order
continuous grass $>$ rotation $>$ continuous corn, with fertilised grass
shifting from a P sink to a source \parencite{Zhang2015b}. Swine-manure
compost produced much greater DRP, PP, and TP than controls
\parencite{Zhang2015}, and about $60\%$ of DRP losses from a no-till
cover-crop site followed fall surface fertilisation, highlighting incidental
loss \parencite{Macrae2023,Mumbi2025b}. Across the MANAGE database, generally
less than $2\%$ of applied P was lost in drainage \parencite{Christianson2016}.
Tillage involves a trade-off between dissolved and particulate forms: no-till
and conservation tillage were associated with higher subsurface DRP, whereas
tillage increased PP and sediment
\parencite{Ni2020,Macrae2023,Koskiaho2002}. In snowmelt regions, zero tillage
and perennial forages can \emph{increase} dissolved P despite reducing erosion
\parencite{Baulch2019}. Overall, effective P control depends on matching the
right practice to site hydrology, connectivity, weather, and legacy P, and on
recognising these dissolved--particulate trade-offs
\parencite{King2015,Ni2020,Baulch2019}.